\subsection{KT-POLY-FAIL: the maximal geometric prime-to-p quotient
can lose a whole fibre}
\label{sec:kt-poly-fail}

\begin{proposition}[KT-POLY-FAIL]
There is a smooth hyperbolic polycurve surface $X/\mathbf F_4$
and two distinct $\mathbf F_4$-points whose section classes in
$\Pi_X^{(2')}$ are equal. Thus the unrestricted extension of
AN\_GRH A05 from curves to hyperbolic polycurves, using the
maximal prime-to-$p$ quotient of the entire geometric fundamental
group, is false.
\end{proposition}

\begin{proof}
Write $k=\mathbf F_4=\mathbf F_2(\omega)$ with
$\omega^2+\omega+1=0$. Let $F/k$ be the smooth projective
genus-two curve
\[
 y^2+y=x^5,
\]
and write $\iota(x,y)=(x,y+1)$ for its hyperelliptic
involution. The unique point $q$ over $x=\infty$ is $k$-rational
and fixed by $\iota$. Set
\[
 B=\mathbf P^1_k\setminus\{0,1,\infty\},\qquad
 B'\longrightarrow B:\quad z^2+z=t+\omega.
                                                               \tag{PF.1}
\]
The second map is a geometrically connected finite \'etale
double cover, with deck involution $\tau(z)=z+1$.
Define
\[
 X=(F\times_k B')/\langle(\iota,\tau)\rangle.
                                                               \tag{PF.2}
\]
The diagonal action is free because its action on $B'$ is free.
The quotient exists (the product is quasi-projective), and finite
\'etale descent shows that $X\to B$ is smooth and proper with
geometric fibre $F$. Together with $B\to\operatorname{Spec}k$
this makes $X$ a hyperbolic polycurve. The fixed point $q$ descends
to a section of $X\to B$.

We work first over $\overline k$. Proper K\"unneth, SGA~1,
Expos\'e X, Corollaire 1.7, pp.~203--204, applied after pullback
to $B'$, and the section supplied by $q$, identify
\[
 \pi_1(X_{\overline k})\simeq
 D\rtimes H,\qquad
 D=\pi_1(F_{\overline k},q),\quad H=\pi_1(B_{\overline k}).
                                                               \tag{PF.3}
\]
The action of $H$ factors through the double-cover character
$H\twoheadrightarrow C_2$ and the induced pointed action
$\iota_*$ on $D$. To spell out the injection implicit in
(PF.3), the subgroup associated with $B'$ is
$D\times\pi_1(B'_{\overline k})$ by proper K\"unneth;
finite \'etale covers inject their fundamental groups. Its
quotient by $D$ is $\pi_1(B'_{\overline k})$, and descent along
$B'/B$ gives $H$. The descended fixed-point section splits this
extension and supplies exactly the stated action.

Let $\varphi:D\rtimes H\to Q$ be any finite quotient of odd
order. Pick $h\in H$ with nontrivial image in $C_2$.
On $\varphi(D)$, conjugation by $\varphi(h)$ is the involution
induced by $\iota_*$. Its order divides two. Its order also
divides the odd order of $\varphi(h)$, so it is the identity.
Hence the restriction $D\to\varphi(D)$ identifies $d$ and
$\iota_*(d)$ for every $d\in D$.

We claim that every finite odd quotient $D\twoheadrightarrow A$
on which $\iota_*$ is trivial is trivial. Let $Y\to F_{\overline k}$
be its connected \'etale Galois cover. The pointed invariance
lifts $\iota$ to an involution $j$ of $Y$ fixing the chosen point
over $q$ and commuting with $A$. Consequently
$Y\to\mathbf P^1_{\overline k}$ is Galois with group
$A\times\langle j\rangle$. Since $Y\to F_{\overline k}$ is
\'etale, every nontrivial inertia group of this composite has
order two. A direct product of an odd group with $C_2$ has a
unique element of order two, namely $j$. Thus the quotient
$Y/\langle j\rangle\to\mathbf P^1_{\overline k}$ is \'etale
with Galois group $A$. The geometric projective line has no
nontrivial finite \'etale connected cover (Riemann--Hurwitz),
so $A=1$. Applying this to $\varphi(D)$ proves
\[
 D\longrightarrow\pi_1(X_{\overline k})^{(2')}
 \quad\hbox{is trivial}.                                  \tag{PF.4}
\]
In fact (PF.3) now gives
$\pi_1(X_{\overline k})^{(2')}\simeq H^{(2')}$.

At $b=\omega\in B(k)$ the double cover (PF.1) has the two
rational points $z=0,1$. Choosing $z=0$ identifies the fibre
$X_b$ with $F$ over $k$. The points $(x,y)=(0,0)$ and $(0,1)$
of $F(k)$ give two distinct points $P,Q\in X_b(k)$.
The map $\pi_1(X_b)\to\Pi_X^{(2')}$ kills its geometric
subgroup by (PF.4). It therefore factors through the quotient
$G_k$, and both point sections of $\pi_1(X_b)\to G_k$ have
the same image. Their sections in $\Pi_X^{(2')}$ are equal
up to the usual geometric conjugacy. This proves the
counterexample.
\end{proof}

\paragraph{The precise obstruction.}
The missing operation is preservation of the fibre group when
passing from a hyperbolic fibration to the maximal prime-to-$p$
quotient of the entire geometric group. Here geometric monodromy
has order $p=2$. It disappears in every prime-to-$p$ quotient,
forcing the fibre's hyperelliptic involution to act trivially and
thereby killing the fibre. A fibrewise prime-to-$p$ completion
that retains the full geometric base group is a different group;
KT-POLY-FAIL makes no assertion against that alternative.

